\section{Used tools}


\subsection{Agno Framework}

Agno is an open-source, Python-based agentic framework designed to facilitate the development of intelligent AI agents with integrated memory, knowledge, tools, and reasoning capabilities. Agno provides a lightweight, composable architecture that emphasizes declarative agent construction, model-agnostic design, and first-class tool support. The framework enables rapid prototyping and deployment of multi-agent systems through a clean API that abstracts the complexities of agent orchestration, tool coordination, and state management while maintaining full transparency for debugging and auditability.


\subsubsection{Tool System and Toolkits}

Tools occupy a central role in Agno's architecture as primary mechanisms for agent interaction with external systems. Tools are Python functions or classes that expose specific capabilities to agents, enabling interaction with external systems such as APIs, databases, web search engines, and file systems. The framework provides two primary tool abstractions: individual \textbf{function tools} for single operations and \textbf{toolkits} for collections of related functions that share internal state and coordinate behavior. Toolkits are implemented as Python classes that register multiple methods as callable tools, allowing agents to invoke them transparently based on task requirements.

\subsubsection{RAG Tools for Information Retrieval}

In the thesis implementation, four core toolkits from Agno's pre-built collection specifically utilized, to enable comprehensive information retrieval across web and academic sources:

\textbf{TavilyTools:} Provides advanced web search capabilities through the Tavily API, a search engine optimized for LLM-driven applications. The toolkit supports configurable search depth (\texttt{standard} vs.\ \texttt{deep}), maximum token limits per result, and output formatting (sourcedAnswer for synthesized responses with citations vs.\ searchResults for raw context data). TavilyTools was employed in the web research branch to execute primary web searches with advanced depth settings, retrieving up to 5 documents per query with full raw content preserved for subsequent LLM-based parsing and summarization.

\textbf{ArxivTools:} Enables retrieval of academic preprints and research papers from the arXiv repository. The toolkit queries the arXiv API using structured search parameters and returns paper metadata including titles, authors, abstracts, publication dates, and arXiv identifiers. ArxivTools was integrated into the academic research branch to complement web-based sources with peer-reviewed and domain-specific scientific literature, ensuring coverage of cutting-edge research findings and formal scholarly discourse.

\textbf{DuckDuckGoTools:} Provides privacy-focused web search through the DuckDuckGo API, serving as a fallback search backend when Tavily or other primary tools are unavailable or rate-limited. The toolkit executes keyword-based searches and returns result snippets, titles, and URLs. While not the primary search tool in the implementation, DuckDuckGoTools was configured as an alternative within the agent toolkit stack to ensure search robustness and redundancy.

\textbf{GoogleSearchTools:} Enables search through Google's Custom Search Engine API, providing access to indexable web content with support for configurable result limits and ranking. The toolkit returns structured results including titles, snippets, and URLs. GoogleSearchTools serves as an alternative search provider within the agent toolkit stack, offering broader web coverage and integration with Google's search index while being subject to rate limits and quota constraints.

\subsubsection{Agent Execution and Tool Coordination}

LangGraph gents execute tools through an internal reasoning loop that analyzes user queries, determines which tools are necessary, invokes them sequentially or in parallel, and synthesizes their outputs into coherent responses. Tool calls are transparent by default (configurable via \texttt{show\_tool\_calls}), enabling inspection of agent decision-making and facilitating debugging. In my implementation, agents were instantiated with the Ollama backend and configured with specific toolkits based on branch requirements (TavilyTools for web research, ArxivTools for academic research).

\subsubsection{Integration of Agno Tools into LangGraph Workflows}

Agno tools integrate seamlessly into LangGraph state machines as callable functions within graph-based workflows. In the developed architecture, each research branch (web and academic) instantiates branch-specific toolkits, executes queries through the tool coordination layer, and returns raw search results to the LangGraph state for subsequent parsing, deduplication, and summarization. This integration pattern leverages Agno's composable tool abstractions while maintaining LangGraph's control over overall workflow orchestration, state management, and conditional routing.




\subsection{DeepEval}

DeepEval is an open-source, LLM-powered evaluation framework designed to assess the quality of RAG systems through quantitative, interpretable metrics. Rather than relying on manual evaluation or subjective criteria, DeepEval employs the LLM-as-a-judge paradigm, where a language model autonomously evaluates generated outputs against multiple complementary dimensions. This systematic approach enables comprehensive quality assurance of generated content without human annotation.

\subsubsection{Core RAG Evaluation Metrics}

DeepEval provides four fundamental built-in metrics specifically designed for RAG pipeline evaluation:

\textbf{Faithfulness:} Measures whether the generated output's claims are factually grounded in the retrieved documents. Faithfulness answers: \textit{What fraction of the generated statements are supported by the retrieval context?} This metric prevents hallucination, the generation of information not present in sources, by validating each factual claim against provided documents. Formally:

\begin{equation}
\text{Faithfulness} = \frac{\text{Supported Claims}}{\text{Total Claims}}
\end{equation}

\textbf{Answer Relevancy:} Evaluates whether the generated output directly addresses the user's original query. Answer relevancy measures topical alignment by classifying generated statements as relevant or tangential to the research question. This metric ensures the output focuses on answering what was asked, independent of factual correctness. Formally:

\begin{equation}
\text{Answer Relevancy} = \frac{\text{Relevant Statements}}{\text{Total Statements}}
\end{equation}

\textbf{Contextual Relevancy:} Assesses the quality of retrieved documents by measuring whether they are genuinely relevant to the user's query. Unlike faithfulness, which validates the generator's use of context, contextual relevancy validates the retriever's selection process. This metric identifies whether the retrieval system has selected appropriate source documents. Formally:

\begin{equation}
\text{Contextual Relevancy} = \frac{\text{Relevant Retrieval Documents}}{\text{Total Retrieved Documents}}
\end{equation}

\textbf{Hallucination Detection:} Explicitly identifies claims in the generated output that contradict or lack support in the retrieval context. This metric flags negative cases, statements that directly contradict sources or introduce unsupported assertions, distinct from positive alignment measured by faithfulness. Formally:

\begin{equation}
\text{Hallucination Score} = 1.0 - \frac{\text{Contradicted/Unsupported Claims}}{\text{Total Claims}}
\end{equation}

\subsubsection{Metric Scoring and Threshold-Based Evaluation}

Each metric produces a score in the range \([0.0, 1.0]\), where 1.0 represents perfect performance and 0.0 represents complete failure. The interpretation of metric scores typically employs a threshold-based approach: scores above a defined threshold are considered acceptable, while scores below the threshold indicate quality concerns requiring investigation or remediation.

The threshold selection is a design choice that reflects the acceptable quality level for the application. In academic or high-stakes applications, thresholds may be set higher (e.g., 0.70, 0.75) to ensure maximum reliability. In exploratory or low-stakes scenarios, thresholds may be lower (e.g., 0.50, 0.55) to permit faster iteration. The threshold value represents a trade-off between precision (few false positives) and recall (few false negatives):

\begin{itemize}
    \item \textbf{High Threshold} (e.g., 0.75): Stringent quality standards; only high-confidence outputs pass; fewer acceptable results but higher confidence in accepted outputs.
    
    \item \textbf{Moderate Threshold} (e.g., 0.50--0.65): Balanced approach; reasonable quality standards; most practical applications employ thresholds in this range.
    
    \item \textbf{Low Threshold} (e.g., 0.30--0.50): Permissive standards; accepts lower-confidence outputs; useful for early-stage prototyping or when perfect accuracy is infeasible.
\end{itemize}


\subsubsection{LLM-as-a-Judge Evaluation Paradigm}

DeepEval's core mechanism relies on the LLM-as-a-judge paradigm: an evaluator language model assesses the quality of generated outputs by applying metric-specific evaluation logic. This approach offers several advantages:

\begin{itemize}
    \item \textbf{Semantic Understanding}: The evaluator LLM comprehends nuanced meaning, context, and domain-specific correctness beyond simple pattern matching.
    
    \item \textbf{Interpretability}: Each metric computation includes reasoning, enabling users to understand why outputs passed or failed evaluations.
    
    \item \textbf{Scalability}: Evaluation can be automated across arbitrary datasets without requiring human annotators, enabling rapid system iteration.
    
    \item \textbf{Consistency}: Uniform evaluation criteria are applied to all outputs, reducing variability inherent in manual assessment.
\end{itemize}

The evaluator operates independently from the generation model, allowing for objective assessment of generation quality without introducing bias from the same model producing the output.

\subsubsection{RAG Pipeline Coverage}

The four metrics comprehensively cover distinct components of the RAG pipeline architecture:

\begin{center}
\begin{tabular}{|l|l|l|}
\hline
\textbf{Pipeline Component} & \textbf{Metric} & \textbf{Assessment Focus} \\
\hline
Retriever Quality & Contextual Relevancy & Document selection appropriateness \\
\hline
Generator Grounding & Faithfulness & Output grounding in sources \\
\hline
Generator Alignment & Answer Relevancy & Query-output topical alignment \\
\hline
Generator Accuracy & Hallucination Detection & Contradiction and fabrication \\
\hline
\end{tabular}
\end{center}

This multi-dimensional evaluation design enables diagnostic analysis of RAG failures. Contextual relevancy failures indicate retriever problems; faithfulness failures indicate generator misuse of context; answer relevancy failures indicate output drift from the query; hallucination failures indicate factual inaccuracies.

\subsubsection{Multi-Metric Evaluation Composition}

DeepEval supports composite evaluation strategies where multiple metrics are evaluated simultaneously on the same output. A typical approach combines all four metrics to form a comprehensive evaluation profile:

\begin{itemize}
    \item Each metric provides independent assessment of a specific quality dimension
    \item Individual metric scores reveal which pipeline components are performing adequately
    \item Combined evaluation results enable holistic quality determination across all dimensions
\end{itemize}

Different applications may emphasize different metrics based on their primary concerns. A system prioritizing factual accuracy would emphasize faithfulness and hallucination detection, while a system prioritizing user satisfaction might emphasize answer relevancy. This flexibility allows to customize evaluation profiles to match application-specific quality objectives.
